\documentclass[a4paper,14pt]{extreport}

\usepackage{graphicx}

% Пакеты для русского языка и кодировки
\usepackage[T2A]{fontenc}
\usepackage[utf8]{inputenc}
\usepackage[russian]{babel}

% Пакеты для математики и теорем
\usepackage{amsmath, amssymb, amsthm}

% Пакеты для графики и гиперссылок
\usepackage{hyperref}

% Настройка полей
\usepackage{geometry}
\geometry{top=2cm, bottom=2cm, left=2.5cm, right=1.5cm}

% Для оформления кода и таблиц
\usepackage{listings}
\usepackage[table]{xcolor}

% Настройка межстрочного интервала
\usepackage{setspace}
\onehalfspacing

% Отключает нумерацию для part, chapter, section и т.д.
\setcounter{secnumdepth}{-3}

\hypersetup{
    colorlinks=true,
    linkcolor=black,
    urlcolor=blue
}


\begin{document}

\hypersetup{pageanchor=false}
\begin{titlepage}
    \centering
    {\large Национальный исследовательский университет «Высшая школа экономики»\\}
    \vspace{2cm}
    {\Large Программа учебной дисциплины\\}
    \vspace{1cm}
    {\LARGE \textbf{Адаптационный курс к дисциплине «Алгоритмы и алгоритмические языки»}\\}
    \vspace{1cm}
    {\large Раздел 2. Основы архитектуры, операционных систем и инструментов программиста\\}
    \vfill
    {Москва\\2026}
\end{titlepage}
\hypersetup{pageanchor=true}

\tableofcontents
\newpage

\chapter{Общие сведения о дисциплине}

\section{Аннотация}

Адаптационный курс предназначен для студентов первого курса образовательной программы «Программная инженерия» и направлен на выравнивание стартовой подготовки перед изучением дисциплины «Алгоритмы и алгоритмические языки». Основное внимание уделяется темам, которые обычно вызывают трудности при переходе от школьной информатики к систематическому программированию: системам счисления, представлению данных в памяти, базовой архитектуре компьютера, роли операционной системы, работе в Linux, сборке программ на C и NASM, а также начальному использованию Git.

Настоящий документ фиксирует место второго раздела курса в общей программе учебной дисциплины. Раздел «Основы архитектуры, операционных систем и инструментов программиста» связывает математические основы представления данных с практической работой программиста: запуском программы, адресами, памятью, процессами, командной строкой, компиляцией, ассемблированием, линковкой и управлением исходным кодом.

\section{Цель освоения дисциплины}

Цель дисциплины состоит в том, чтобы подготовить студентов к дальнейшему изучению программирования на C и NASM, сформировать базовую техническую грамотность и дать практические навыки работы с инструментами, которые будут использоваться в основном курсе.

\section{Место раздела 2 в структуре дисциплины}

Раздел 2 изучается после темы о системах счисления и машинном представлении чисел. Он служит переходом от представления данных к исполнению программ: студент должен увидеть, что двоичные данные, память, процессор, операционная система и инструменты сборки являются частями одной рабочей модели.

Раздел подготавливает студентов к следующим видам деятельности:

\begin{itemize}
    \item написание простых программ на C;
    \item чтение и запуск примеров на NASM;
    \item работа с Linux как основной учебной средой;
    \item диагностика типовых ошибок запуска и сборки;
    \item ведение учебных проектов с помощью Git.
\end{itemize}

\chapter{Планируемые результаты обучения по разделу 2}

После освоения раздела студент должен знать:

\begin{itemize}
    \item назначение основных компонентов компьютера: CPU, RAM, SSD/HDD, GPU;
    \item базовые понятия памяти: бит, байт, машинное слово, адрес, адресное пространство;
    \item роль регистров процессора, стека и кучи на уровне общей модели исполнения программы;
    \item отличие программы как файла от процесса как выполняющегося экземпляра программы;
    \item назначение операционной системы и её роль в управлении процессами, памятью, файлами, устройствами и правами доступа;
    \item назначение стандартных потоков \texttt{stdin}, \texttt{stdout}, \texttt{stderr};
    \item основные этапы сборки C-программы и NASM-программы;
    \item назначение \texttt{gcc}, \texttt{nasm}, \texttt{ld}, \texttt{make}, \texttt{gdb}, \texttt{objdump}, \texttt{readelf}, \texttt{nm};
    \item назначение Git и базовую модель репозитория, индекса, коммита и удалённого репозитория.
\end{itemize}

После освоения раздела студент должен уметь:

\begin{itemize}
    \item объяснять, где хранится программа до запуска и что происходит при её выполнении;
    \item связывать переменные, массивы и указатели в C с адресами и байтами в памяти;
    \item запускать программы из командной строки Linux;
    \item использовать права доступа и исполняемый бит;
    \item работать с каталогами и файлами в терминале;
    \item использовать перенаправление потоков и пайпы;
    \item устанавливать и проверять инструменты разработки в Debian/Linux;
    \item собирать и запускать простые программы на C;
    \item собирать и запускать простые программы на NASM под Linux x86-64;
    \item читать сообщения об ошибках компиляции, ассемблирования, линковки и запуска;
    \item выполнять минимальный цикл работы с Git: \texttt{init}, \texttt{status}, \texttt{add}, \texttt{commit}, \texttt{log}, \texttt{diff}, \texttt{clone}, \texttt{pull}, \texttt{push}.
\end{itemize}

\chapter{Содержание раздела 2}

\section{Тема 3. Основы архитектуры компьютера}

\subsection*{Лекция 4. Как устроен компьютер и где исполняется программа}

Основные компоненты компьютера: CPU, RAM, SSD/HDD, GPU. Бит, байт, машинное слово, адрес, адресное пространство. Регистр процессора как самая быстрая рабочая память CPU. Оперативная память как массив байтов с адресами. Отличие хранения программы на диске от исполнения программы в памяти. Стек и куча на уровне общей идеи. Связь с C: переменная, указатель, массив, размер типа. Связь с NASM: регистры, инструкции, обращение к памяти.

\subsection*{Семинар 4. Память, адреса и представление данных}

Практическая работа с размерами базовых типов C через \texttt{sizeof}. Представление массивов \texttt{char}, \texttt{int}, \texttt{long} в памяти. Вычисление адреса элемента массива по адресу начала массива и размеру элемента. Различие между значением переменной и адресом переменной. Просмотр байтов целого числа в памяти и объяснение порядка байт. Начальное знакомство с \texttt{gdb} и командами просмотра памяти.

\section{Тема 4. Операционная система и запуск программы}

\subsection*{Лекция 5. Операционная система и запуск программы}

Задачи операционной системы: процессы, память, файлы, устройства, права доступа. Программа как файл на диске и процесс как выполняющийся экземпляр программы. Загрузка программы в память и передача управления начальной точке исполнения. Аргументы командной строки и переменные окружения. Стандартные потоки \texttt{stdin}, \texttt{stdout}, \texttt{stderr}. Код завершения программы. Права доступа к файлам и исполняемый бит. Отличия Windows, Linux и macOS с точки зрения программиста. Формат ELF, системные вызовы, пользовательский режим и режим ядра на уровне общей идеи.

\subsection*{Семинар 5. Процессы, файлы и стандартные потоки в Linux}

Запуск готовой программы из текущего каталога через \texttt{./program}. Проверка типа файла командой \texttt{file}. Изменение прав доступа через \texttt{chmod +x}. Передача аргументов командной строки. Перенаправление вывода в файл и входа из файла. Соединение команд через пайп. Проверка кода завершения через \texttt{echo \$?}. Диагностика ошибок \texttt{permission denied}, \texttt{command not found}, \texttt{no such file or directory}.

\section{Тема 5. Linux как рабочая среда программиста}

\subsection*{Лекция 6. Терминал, shell и файловая система Linux}

Гипервизор, виртуальная машина, WSL2 и Docker как способы получить изолированную рабочую среду. Debian/Linux как базовая учебная среда для курса. Разница между установкой Debian второй системой и работой в виртуальной машине: преимущества, ограничения, риски и разумные сценарии выбора. Терминал и shell: кто читает команду пользователя. Файловая система Linux: корень \texttt{/}, домашний каталог, текущий каталог. Абсолютные и относительные пути. Каталоги \texttt{/home}, \texttt{/bin}, \texttt{/usr}, \texttt{/tmp} на уровне пользователя. Переменная окружения \texttt{PATH}. Работа с документацией через \texttt{man} и \texttt{--help}.

\subsection*{Семинар 6. Навигация, поиск и работа с текстом}

Отработка минимального набора команд Linux: \texttt{pwd}, \texttt{cd}, \texttt{ls}, \texttt{mkdir}, \texttt{touch}, \texttt{cp}, \texttt{mv}, \texttt{rm}, \texttt{less}, \texttt{grep}, \texttt{find}, \texttt{wc}. Создание рабочего каталога для курса. Поиск файлов с расширениями \texttt{.c} и \texttt{.asm}. Просмотр текстовых файлов и результатов работы программ. Построение простых пайплайнов. Использование базовых glob-паттернов и регулярных выражений в \texttt{grep}.

\section{Тема 6. Сборка программ и инструменты разработки}

\subsection*{Лекция 7. От исходного кода к исполняемому файлу}

Исходный файл, объектный файл, исполняемый файл. Этапы сборки C-программы: препроцессор, компилятор, ассемблер, линковщик. Заголовочные файлы и единицы трансляции. Что делает \texttt{gcc} при обычном запуске. Что делает \texttt{nasm} при сборке ассемблерного файла. Зачем нужен линковщик. Зачем нужен \texttt{make} и файл \texttt{Makefile}. Типовые ошибки: ошибка компиляции, ошибка ассемблирования, ошибка линковки. Секции исполняемого файла \texttt{.text}, \texttt{.data}, \texttt{.bss}; символы и таблица символов; ABI и соглашение о вызовах на уровне подготовки к дальнейшему курсу.

\subsection*{Семинар 7. Сборка C- и NASM-программ}

Сборка программы \texttt{hello.c} командой \texttt{gcc hello.c -o hello}. Запуск получившегося исполняемого файла. Разделение C-программы на несколько исходных файлов и сборка их вместе. Получение объектного файла через \texttt{gcc -c}. Написание минимального \texttt{Makefile} с целями \texttt{all} и \texttt{clean}. Сборка NASM-программы в объектный файл командой \texttt{nasm -f elf64}. Линковка объектного файла в исполняемый файл. Просмотр секций через \texttt{readelf -S}, дизассемблирование через \texttt{objdump -d}, поиск символов через \texttt{nm}.

\section{Тема 7. Основы Git и командной разработки}

\subsection*{Лекция 8. Система контроля версий Git}

Назначение системы контроля версий. Репозиторий, рабочая директория, индекс, коммит. История изменений и идентификатор коммита. Локальный и удалённый репозиторий. Ветка как линия разработки. Конфликт и причины его возникновения. Минимальная дисциплина работы: маленькие коммиты, понятные сообщения, проверка изменений перед коммитом, исключение артефактов сборки.

\subsection*{Семинар 8. Минимальный рабочий цикл Git}

Создание локального репозитория командой \texttt{git init}. Проверка состояния рабочей директории через \texttt{git status}. Добавление файла в индекс через \texttt{git add}. Создание коммита через \texttt{git commit}. Просмотр истории через \texttt{git log}. Просмотр изменений через \texttt{git diff}. Клонирование удалённого репозитория через \texttt{git clone}. Получение и отправка изменений через \texttt{git pull} и \texttt{git push}. Создание ветки, переключение между ветками и учебное разрешение конфликта.

\chapter{Элементы контроля}

Для раздела 2 используются следующие элементы контроля:

\begin{itemize}
    \item аудиторная работа на семинарах: выполнение практических заданий, ответы на вопросы преподавателя, участие в разборе ошибок;
    \item практическое задание по архитектуре и памяти: объяснение размера типов C, адресации массива и просмотра байтов значения в памяти;
    \item практическое задание по Linux: навигация по каталогам, работа с файлами, перенаправление потоков, пайпы, запуск программ;
    \item практическое задание по сборке: сборка и запуск простой программы на C и простой программы на NASM;
    \item практическое задание по Git: создание репозитория, коммит изменений, просмотр истории и различий;
    \item мини-контрольная или устный опрос по базовым понятиям: память, адрес, процесс, исполняемый файл, стандартные потоки, код завершения, этапы сборки.
\end{itemize}

\chapter{Промежуточная аттестация}

Формула оценки результатов освоения раздела 2:

\begin{itemize}
    \item 40\% -- аудиторная работа и выполнение заданий на семинарах;
    \item 30\% -- практическое задание по Linux и запуску программ;
    \item 20\% -- практическое задание по сборке C/NASM;
    \item 10\% -- задание по Git или короткий теоретический опрос.
\end{itemize}

Блокирующим элементом является практическая проверка рабочей среды: студент должен показать, что умеет открыть терминал, перейти в рабочий каталог, собрать и запустить C-программу, собрать и запустить NASM-программу, проверить код завершения и объяснить одну типовую ошибку сборки или запуска.

\chapter{Список литературы и ресурсов}

\begin{itemize}
    \item Bryant R. E., O'Hallaron D. R. Computer Systems: A Programmer's Perspective.
    \item Kernighan B. W., Ritchie D. M. The C Programming Language.
    \item Patterson D. A., Hennessy J. L. Computer Organization and Design.
    \item Debian Documentation: \url{https://www.debian.org/doc/}
    \item The Linux Command Line by William Shotts: \url{https://linuxcommand.org/tlcl.php}
    \item GNU Make Manual: \url{https://www.gnu.org/software/make/manual/}
    \item NASM Documentation: \url{https://www.nasm.us/docs.php}
    \item Git Book: \url{https://git-scm.com/book/}
    \item GDB Documentation: \url{https://sourceware.org/gdb/documentation/}
\end{itemize}

\chapter{Авторы и разработчики}

Программа раздела определяет содержание, результаты обучения и элементы контроля в рамках адаптационного курса факультета Программной инженерии НИУ ВШЭ.

\end{document}
